% Compilar con LuaLaTeX o XeLaTeX (Overleaf: Menu > Compiler > LuaLaTeX)
\documentclass[11pt,a4paper]{article}
\usepackage{fontspec}
\usepackage[spanish]{babel}
\usepackage[a4paper,margin=2.5cm,headheight=15pt]{geometry}
\usepackage{graphicx,xcolor,float,booktabs,array,tabularx,fancyhdr,titlesec,enumitem,caption,fvextra}
\usepackage[hidelinks]{hyperref}
\definecolor{ujablue}{HTML}{1F4E8C}
\definecolor{ujapink}{HTML}{C2185B}
\definecolor{gris}{HTML}{666666}
\setlength{\parindent}{0pt}\setlength{\parskip}{0.6em}
\pagestyle{fancy}\fancyhf{}
\lhead{\small\color{gris}Gestión de Proyectos Software --- UJA}
\rhead{\small\color{gris}Práctica 3: Arquitectura Ágil y UML con IA}
\lfoot{\small\color{gris}Grupo \textbf{G11}}\rfoot{\small\color{gris}\thepage}
\renewcommand{\headrulewidth}{0.4pt}\renewcommand{\footrulewidth}{0.4pt}
\titleformat{\section}{\Large\bfseries\color{ujablue}}{\thesection.}{0.6em}{}
\titleformat{\subsection}{\large\bfseries\color{ujablue}}{\thesubsection.}{0.6em}{}
\titleformat{\paragraph}[runin]{\bfseries}{}{0pt}{}[.]
\fvset{fontsize=\scriptsize,breaklines=true,frame=single,framesep=2mm,rulecolor=\color{gris!50}}
\captionsetup{font=small}

% \bloqueGlobal{clave}{titulo}{alcance}   (Estructura A)
\newcommand{\bloqueGlobal}[3]{%
 \subsection{#2}
 \paragraph{Alcance y contexto de partida} #3
 \paragraph{Prompt enviado a la IA}\par\VerbatimInput{prompts/#1.txt}
 \paragraph{Auditoría e informe de errores}\par\input{auditorias/#1.tex}
 \paragraph{Código Mermaid corregido y diagrama renderizado}\par\VerbatimInput{codigo/#1.mmd}
 \begin{figure}[H]\centering
 \IfFileExists{figuras/#1.png}{\includegraphics[width=\linewidth,height=0.45\textheight,keepaspectratio]{figuras/#1.png}}{\fbox{\textcolor{red}{Falta figuras/#1.png}}}
 \caption{#2.}\end{figure}}

% \bloqueHU{clave}{id}{titulo}{texto HU}{criterios}   (Estructura B)
\newcommand{\bloqueHU}[5]{%
 \subsection{#2. #3}
 \paragraph{Requisito / HU de partida} \textbf{#2.} #4 \emph{Criterios de aceptación:} #5
 \paragraph{Prompt enviado a la IA}\par\VerbatimInput{prompts/#1.txt}
 \paragraph{Auditoría e informe de errores}\par\input{auditorias/#1.tex}
 \paragraph{Código Mermaid corregido y diagrama renderizado}\par\VerbatimInput{codigo/#1.mmd}
 \begin{figure}[H]\centering
 \IfFileExists{figuras/#1.png}{\includegraphics[width=\linewidth,height=0.4\textheight,keepaspectratio]{figuras/#1.png}}{\fbox{\textcolor{red}{Falta figuras/#1.png}}}
 \caption{#2: #3.}\end{figure}}

\begin{document}
\begin{titlepage}
\thispagestyle{empty}\centering
\vspace*{1cm}
\IfFileExists{logo_uja.png}{\includegraphics[height=3cm]{logo_uja.png}}{\fbox{\textcolor{red}{Colocar logo\_uja.png}}}\\[1.5cm]
{\Large\bfseries\color{ujablue}UNIVERSIDAD DE JAÉN}\\[0.3cm]
Gestión de Proyectos Software\\[1.2cm]
{\color{ujablue}\rule{\linewidth}{0.8pt}}\\[1cm]
{\fontsize{34}{40}\selectfont\bfseries\color{ujablue}Práctica 3}\\[0.4cm]
{\Large\bfseries Arquitectura Ágil: Diagramación UML Completa\\ Asistida por IA y Mermaid}\\[0.8cm]
Proyecto \textbf{Atelier Jaén}: comercio electrónico para una boutique de moda femenina\\[0.6cm]
{\color{ujapink}\rule{0.45\linewidth}{0.8pt}}\\[1cm]
{\large\bfseries Integrantes del grupo G11}\\[0.4cm]
\begin{tabular}{@{}lll@{}}
\toprule
\textbf{Nombre y apellidos} & \textbf{Rol (Práctica 1)} & \textbf{Correo electrónico}\\
\midrule
Juan Luis Fernández Bayona & Product Owner & \texttt{jlfb0001@red.ujaen.es}\\
David Molina Gámez & Scrum Master & \texttt{dmg00068@red.ujaen.es}\\
Antonio José García Arias & Desarrollador / Diseñador & \texttt{ajga0011@red.ujaen.es}\\
\bottomrule
\end{tabular}
\vfill
Grupo G11\\ Curso académico 2026--2027\\ Noviembre de 2026
\end{titlepage}

\tableofcontents
\clearpage
\section{Enlaces de acceso al proyecto}
URL del tablero Miro (Grupo G11): \url{https://miro.com/app/board/uXjVHnAiHj0=/?share_link_id=353339930990}

\section{Contexto y alcance de la práctica}
Esta memoria evoluciona el Product Backlog de la Práctica~1 (diez historias de usuario repartidas en tres sprints) hacia una arquitectura técnica modelada con Mermaid y generada con apoyo de IA. Cada diagrama se documenta con su flujo de auditoría: partida, prompt literal, informe de errores y código corregido con su imagen renderizada.

\begin{center}\small
\begin{tabular}{@{}llp{9cm}@{}}
\toprule \textbf{ID} & \textbf{Sprint} & \textbf{Historia de usuario}\\ \midrule
US-1 & 1 & Visualización del Catálogo Principal\\
US-2 & 1 & Ficha de Detalle de Producto\\
US-3 & 1 & Cesta de la Compra\\
US-4 & 1 & Procesamiento de Pago con Tarjeta (Checkout)\\
US-5 & 1 & Panel Básico de Gestión de Pedidos (PYME)\\
US-6 & 2 & Filtrado por Categoría y Tallas\\
US-7 & 2 & Pago Rápido con Bizum\\
US-8 & 2 & Registro e Inicio de Sesión de Usuario\\
US-9 & 3 & Aplicación de Cupones de Descuento\\
US-10 & 3 & Guía Interactiva de Tallas\\
\bottomrule\end{tabular}
\end{center}

\textbf{Herramientas:} Mermaid (mermaid.live) para el renderizado y [indicar IA utilizada, p.~ej. Claude] para la generación inicial.

\section{Bloque 1: Diagramas de casos de uso}

\bloqueGlobal{cu_general}{Diagrama general de casos de uso (usecase-beta)}{Cubre la totalidad de los Sprints 1, 2 y 3: cuatro paquetes (catálogo, compra y pago, cuenta, administración), cinco actores y las relaciones include/extend entre casos de uso.}
\clearpage
\bloqueHU{cu_us01}{US-1}{Visualización del Catálogo Principal}{Como cliente, quiero ver el catálogo general de prendas y accesorios para conocer los productos disponibles de la tienda.}{Cada producto debe mostrar imagen, nombre, precio con IVA y distintivo de disponibilidad en stock.}
\bloqueHU{cu_us02}{US-2}{Ficha de Detalle de Producto}{Como cliente, quiero acceder al detalle del artículo y seleccionar mi talla para comprobar sus características y elegir la prenda adecuada.}{Debe existir un selector interactivo de talla (S, M, L, etc.) que deshabilite las opciones agotadas y permita añadir el producto al carrito.}
\bloqueHU{cu_us03}{US-3}{Cesta de la Compra}{Como cliente, quiero gestionar los artículos añadidos al carrito para revisar mi pedido antes de tramitarlo.}{El usuario puede modificar cantidades, eliminar artículos y visualizar en tiempo real el desglose de costes y gastos de envío.}
\bloqueHU{cu_us04}{US-4}{Procesamiento de Pago con Tarjeta (Checkout)}{Como cliente, quiero pagar mi pedido mediante tarjeta de débito o crédito para formalizar la compra de manera inmediata.}{Integración con pasarela bancaria segura con soporte 3D Secure. Confirmación en pantalla y envío de email con número de pedido.}
\bloqueHU{cu_us05}{US-5}{Panel Básico de Gestión de Pedidos (PYME)}{Como administradora de la boutique, quiero ver los pedidos entrantes pagados para preparar y gestionar el empaquetado de envíos.}{Tabla privada con lista de pedidos, datos de entrega del cliente, estado (Pendiente/Enviado) y botón para marcar completado.}
\bloqueHU{cu_us06}{US-6}{Filtrado por Categoría y Tallas}{Como clienta, quiero filtrar el catálogo por categorías (Vestidos, Blusas, Pantalones) y por mi talla para encontrar rápidamente lo que busco.}{Filtro dinámico sin necesidad de recargar la página por completo; muestra contador de resultados coincidentes.}
\bloqueHU{cu_us07}{US-7}{Pago Rápido con Bizum}{Como clienta habitual de compras móviles, quiero pagar mediante Bizum para finalizar mi compra sin teclear los números de la tarjeta.}{Redirección a la pasarela oficial de Bizum solicitando el teléfono y autorización en la app del banco del comprador.}
\bloqueHU{cu_us08}{US-8}{Registro e Inicio de Sesión de Usuario}{Como clienta frecuente, quiero registrarme con mi email para guardar mis direcciones y consultar el historial de compras pasadas.}{Formulario con validación de contraseña segura, opción de recordar sesión y panel con listado de pedidos anteriores.}
\bloqueHU{cu_us09}{US-9}{Aplicación de Cupones de Descuento}{Como cliente, quiero ingresar un código promocional en el carrito para obtener un descuento en mi compra.}{Validación contra lista de cupones válidos en la base de datos y deducción visual inmediata del importe en el total.}
\bloqueHU{cu_us10}{US-10}{Guía Interactiva de Tallas}{Como cliente indecisa, quiero consultar una tabla interactiva de medidas anatómicas en centímetros para acertar con mi talla.}{Ventana modal en la ficha del producto con medidas de contorno de pecho, cintura y cadera por cada talla.}
\section{Bloque 2: Modelado estructural de datos}
\bloqueGlobal{er}{Diagrama Entidad-Relación global}{Modelo conceptual de datos que da soporte a las diez historias de los Sprints 1, 2 y 3.}
\bloqueGlobal{clases}{Diagrama de clases UML global}{Especificación técnica de clases, métodos, atributos, visibilidad y relaciones (asociación, agregación, composición y herencia).}
\section{Bloque 3: Diagramas de secuencia por historia de usuario}
\bloqueHU{sec_us01}{US-1}{Visualización del Catálogo Principal}{Como cliente, quiero ver el catálogo general de prendas y accesorios para conocer los productos disponibles de la tienda.}{Cada producto debe mostrar imagen, nombre, precio con IVA y distintivo de disponibilidad en stock.}
\bloqueHU{sec_us02}{US-2}{Ficha de Detalle de Producto}{Como cliente, quiero acceder al detalle del artículo y seleccionar mi talla para comprobar sus características y elegir la prenda adecuada.}{Debe existir un selector interactivo de talla (S, M, L, etc.) que deshabilite las opciones agotadas y permita añadir el producto al carrito.}
\bloqueHU{sec_us03}{US-3}{Cesta de la Compra}{Como cliente, quiero gestionar los artículos añadidos al carrito para revisar mi pedido antes de tramitarlo.}{El usuario puede modificar cantidades, eliminar artículos y visualizar en tiempo real el desglose de costes y gastos de envío.}
\bloqueHU{sec_us04}{US-4}{Procesamiento de Pago con Tarjeta (Checkout)}{Como cliente, quiero pagar mi pedido mediante tarjeta de débito o crédito para formalizar la compra de manera inmediata.}{Integración con pasarela bancaria segura con soporte 3D Secure. Confirmación en pantalla y envío de email con número de pedido.}
\bloqueHU{sec_us05}{US-5}{Panel Básico de Gestión de Pedidos (PYME)}{Como administradora de la boutique, quiero ver los pedidos entrantes pagados para preparar y gestionar el empaquetado de envíos.}{Tabla privada con lista de pedidos, datos de entrega del cliente, estado (Pendiente/Enviado) y botón para marcar completado.}
\bloqueHU{sec_us06}{US-6}{Filtrado por Categoría y Tallas}{Como clienta, quiero filtrar el catálogo por categorías (Vestidos, Blusas, Pantalones) y por mi talla para encontrar rápidamente lo que busco.}{Filtro dinámico sin necesidad de recargar la página por completo; muestra contador de resultados coincidentes.}
\bloqueHU{sec_us07}{US-7}{Pago Rápido con Bizum}{Como clienta habitual de compras móviles, quiero pagar mediante Bizum para finalizar mi compra sin teclear los números de la tarjeta.}{Redirección a la pasarela oficial de Bizum solicitando el teléfono y autorización en la app del banco del comprador.}
\bloqueHU{sec_us08}{US-8}{Registro e Inicio de Sesión de Usuario}{Como clienta frecuente, quiero registrarme con mi email para guardar mis direcciones y consultar el historial de compras pasadas.}{Formulario con validación de contraseña segura, opción de recordar sesión y panel con listado de pedidos anteriores.}
\bloqueHU{sec_us09}{US-9}{Aplicación de Cupones de Descuento}{Como cliente, quiero ingresar un código promocional en el carrito para obtener un descuento en mi compra.}{Validación contra lista de cupones válidos en la base de datos y deducción visual inmediata del importe en el total.}
\bloqueHU{sec_us10}{US-10}{Guía Interactiva de Tallas}{Como cliente indecisa, quiero consultar una tabla interactiva de medidas anatómicas en centímetros para acertar con mi talla.}{Ventana modal en la ficha del producto con medidas de contorno de pecho, cintura y cadera por cada talla.}
\end{document}
